\section{Introduction}
\label{sec:intro}

A transition into an absorbing state eliminates fluctuations. With two
symmetry-related absorbing configurations, it can also select one of two
ordered states. In two dimensions this combination is closely associated
with generalized voter (GV) criticality, whose interface density vanishes
logarithmically~\cite{dornic2001,clifford1973,holley1975,cox1986,krapivsky1992},
or with separate Ising symmetry breaking and directed-percolation
absorption~\cite{droz2003,alhammal2005,vazquez2008,azizi2020}.
A magnetization jump occurs even at a GV transition. Establishing a
discontinuous loss of activity therefore requires a local activity
measurement and evidence about phase competition.

The Langevin description of Al Hammal \emph{et al.} includes a bistable
deterministic drift, but their two-dimensional simulations gave a GV
transition, with preliminary first-order behavior reported in three
dimensions~\cite{alhammal2005}. Subsequent microscopic studies found GV
and Ising--DP scenarios~\cite{castellano2009,azizi2020}.
Rodrigues \emph{et al.} reported a discontinuous paramagnetic--absorbing
line on a triangular lattice using magnetization and susceptibility
measurements~\cite{rodrigues2015}. The question addressed here is narrower
than the first occurrence of a discontinuous line: can finite local activity,
coherent interfacial competition, and nucleation be measured together in a
two-dimensional binary model with exactly two absorbing configurations?

Discontinuous absorbing transitions are established in reaction
models~\cite{ziff1986,bidaux1989,lubeck2006}, in models with three symmetric
absorbing states~\cite{lipowski2002,azizi2018}, and in ecological models with
two competing absorbing alliances~\cite{szabo2004,guisoni2013}.
The alliances contain many configurations and retain internal dynamics;
their fronts and active boundary layers nevertheless provide close
precedents. Our construction uses binary sites and a flip table depending
only on the number of disagreeing neighbors. It belongs to the nonlinear
voter family~\cite{cox1991,molofsky1999,schweitzer2009}.
A positive voter component removes all frozen configurations except the
two uniform states.

We concentrate on $p=\alpha=1/2$. Bulk simulations resolve a dense active
branch across the measured reversal region of independently prepared
fronts. Decay below that region develops absorbing domains after a
long-lived active stage; a disk and a hole exchange their growth
directions on opposite sides. Front measurements at several orientations
test the possibility of orientation-dependent coexistence, which is allowed
out of equilibrium~\cite{wang2015,liu2018}, without assuming a unique
thermodynamic coexistence point. A wall between the two absorbing states
also develops a layer of active material.

The approach to the voter point is a separate question. We examine longer
runs that distinguish a persistent branch from a slowly decaying apparent
plateau. A finite-time voter closure, motivated by perturbation theory
in $d\ge3$~\cite{cox2013}, yields an invasion scale proportional to
$\alpha(1-p)$. Contributions with at least five distinct ancestral
lineages can generate bistability in that closure; matching it to the
two-dimensional stochastic process remains uncontrolled.
The five-cell rule of Molofsky \emph{et al.} supplies a useful comparison,
reported in the supplement: a small-seed growth crossover and the onset
of coherent expansion can be measured there, but a metastable retreating
active slab is not resolved. This distinction limits what the closure
can be said to explain.
